% ---------------------------------------------------------------------
%  第 8 章 移动窗传播
% ---------------------------------------------------------------------
\chapter{移动窗传播}
\label{ch:moving}

完整的时间推进需要保留全部历史：第 $n$ 步要访问 $G(n,j)$，
$j=0,\dots,n$。存储与单步代价都随时间线性增长，总代价为 $O(n^2)$。
当系统在有限时间内趋于稳态或弛豫时，Green 函数对远端时间的依赖
逐渐衰减，于是可以只保留最近 $t_c+1$ 个时间步。本章描述这一近似
在库中的完整实现：容器、存储、推进、初始化、外推、求解器与
观测量。

\section{容器族}
\label{sec:moving-types}

\cref{tab:moving-fields}列出移动窗容器的字段及其职责。

移动窗只保留推迟与小于分量，Matsubara 与左混合分量不在窗口中。
三个容器承担不同角色：
\begin{itemize}
    \item \code{MovingGreen}：移动窗对象，持有两个环形缓冲
          \code{ret} 与 \code{less}，以及领头时间步 \code{tstp}、
          初始化标志与代数计数；
    \item \code{MovingGreenTimeStep}：窗口内单个时间步的视图，
          持有两个 \code{RowComp}；
    \item \code{MovingTimeFunction}：时间相关的单粒子函数
          （例如 $\bm{H}(t)$）的环形缓冲，布局与 Green 函数一致。
\end{itemize}

\begin{table}[htbp]
    \centering
    \caption[移动窗容器的字段]{移动窗容器的字段}
    \label{tab:moving-fields}
    \small
    \begin{tabularx}{\textwidth}{@{}l>{\raggedright\arraybackslash}p{.31\textwidth}>{\raggedright\arraybackslash}X@{}}
        \toprule
        类型 & 字段 & 含义 \\
        \midrule
        \code{MovingGreen} &
            \code{tc}、\code{tstp}、\code{initialized}、\code{sign}、
            \code{generation}、\code{ret}、\code{less} &
            窗口长度、领头时间步、初始化标志、统计符号、代际计数和两个环形缓冲 \\
        \addlinespace
        \code{RingComp} &
            \code{tc}、\code{size1}、\code{size2}、\code{head}、\code{data} &
            环形缓冲：\code{head} 是逻辑行 0 的物理行号 \\
        \addlinespace
        \code{MovingGreenTimeStep} &
            \code{tc}、\code{tstp}、\code{sign}、\code{ret}、\code{less} &
            单个时间步的两个行分量 \\
        \addlinespace
        \code{RowComp} &
            \code{tc}、\code{size1}、\code{size2}、\code{data} &
            一行 \code{tc+1} 个块 \\
        \addlinespace
        \code{MovingTimeFunction} &
            \code{tc}、\code{tstp}、\code{initialized}、\code{head}、
            \code{size1}、\code{size2}、\code{data} &
            时间函数的环形缓冲 \\
        \bottomrule
    \end{tabularx}
\end{table}

\section{环形行存储}
\label{sec:ring-storage}

\cref{fig:ring-buffer}以五个物理槽为例，标出各槽对应的逻辑行。

设窗口的领头时间为 $t_0$，窗口长度为 $t_c$。块
$(\delti,\deltj)$，$0\le\delti,\deltj\le t_c$，保存
\begin{equation}
    \Gret(t_0-\delti h,\; t_0-\delti h-\deltj h),
    \qquad
    \Gless(t_0-\delti h,\; t_0-\delti h-\deltj h).
    \label{eq:moving-block}
\end{equation}
每个分量是一块连续的数组，大小为
\begin{equation}
    (t_c+1)\times(t_c+1)\times s_1\times s_2 .
\end{equation}
逻辑行 $\delti$ 存放在物理行
\begin{equation}
    \code{physical}(\delti) = (\code{head}+\delti) \bmod (t_c+1).
    \label{eq:ring-index}
\end{equation}

环形存储的关键是“逻辑行”与“物理行”的分离。推进一格时，只需把
\code{head} 减一并对 $t_c+1$ 取非负余数，然后清零新的行。推进无需数据
搬移，也没有任何分配。代价是每次访问都要做一次取模，但这个取模
只在行级别发生一次，行内的块仍然连续。文本格式写出的是逻辑行，
而不是物理行，因此同一份数据在不同 \code{head} 下完全等价。
块 $(\delti,\deltj)$ 的偏移为
\begin{equation}
    \big[\code{physical}(\delti)\cdot(t_c+1)+\deltj\big]\cdot s_1 s_2 .
\end{equation}

\begin{figure}[htbp]
    \centering
    \begin{tikzpicture}[scale=0.95]
        % physical rows
        \foreach \r in {0,...,4} {
            \fill[kbblue!12] (0,-\r*0.8) rectangle ++(4.6,0.66);
            \node[font=\scriptsize, anchor=east] at (-0.1,-\r*0.8+0.33)
                {phys \r};
        }
        % highlight the head row (physical 3)
        \draw[very thick, kborange] (-0.03,-2.43) rectangle (4.63,-1.77);
        % logical labels, aligned with row centres
        \node[font=\scriptsize, anchor=west] at (4.8,0.33)  {$\delti=2$};
        \node[font=\scriptsize, anchor=west] at (4.8,-0.47) {$\delti=3$};
        \node[font=\scriptsize, anchor=west] at (4.8,-1.27) {$\delti=4$};
        \node[font=\scriptsize, anchor=west] at (4.8,-2.07) {$\delti=0$};
        \node[font=\scriptsize, anchor=west] at (4.8,-2.87) {$\delti=1$};
        % head marker
        \node[font=\small, black, anchor=north, align=center] at (2.3,-3.75)
            {橙色框为 \code{head}=3，即逻辑行 $0$ 的物理位置};
    \end{tikzpicture}
    \caption[环形行存储]{环形行存储。逻辑行与物理行相差一个循环平移；
             窗口推进时只需修改 \code{head} 并清零新行，不搬移数据。}
    \label{fig:ring-buffer}
\end{figure}

\section{相对索引}
\label{sec:relative-index}

移动窗对外提供绝对时间与相对时间两种访问方式：
\begin{itemize}
    \item \code{ret_block(n, m)}：按绝对时间步 $n,m$ 访问，内部
          换算成 $\delti = t_0-n$、$\deltj = n-m$；
    \item \code{mrow(delti)} / \code{mblock(delti, deltj)}：按相对
          指标访问，不做换算。
\end{itemize}
换算的前提是两个时间步都落在窗口内：
\begin{equation}
    0 \le t_0-n \le t_c,\qquad
    0 \le n-m \le t_c .
\end{equation}
越界访问立即 panic。这一检查在热循环中会被编译器部分消除：
内核通常先用相对指标确认窗口覆盖，再进入循环。

\section{行旋转}
\label{sec:rotation}

窗口向前推进一格时，最旧的一行被丢弃，新的领头行被写入。实现
不搬移数据，只做两次操作：
\begin{enumerate}[label=(\arabic*)]
    \item 更新头指针
          \begin{equation}
              \code{head} \leftarrow
              (\code{head} + t_c) \bmod (t_c+1),
          \end{equation}
          即向前回退一格；
    \item 把新的领头行清零。
\end{enumerate}
随后求解器把本步结果写入逻辑行 0。\code{generation} 计数在每次
旋转时加一，用于检测过期的视图。

\section{初始化}
\label{sec:moving-init}

移动窗有四个初始化入口：
\begin{itemize}
    \item \code{init_from_green}：由完整 Green 函数及其伴随伙伴
          填充窗口。窗口覆盖 $t_0-t_c,\dots,t_0$，若 $t_c$ 小于
          $t_0$，更早的历史被丢弃；
    \item \code{init_from_green_herm}：在满足厄米关系时只需一个
          完整对象，另一半由共轭重构；
    \item \code{init_from_function}：由时间函数初始化
          \code{MovingTimeFunction}；
    \item \code{init_from_steady}：由稳态对象按相对时间填充，
          用于长时间驱动的稳态初值。
\end{itemize}
初始化完成后 \code{initialized} 置为真。未初始化的对象在访问数据
或推进时 panic，而不是返回零。

\section{时间步视图}
\label{sec:moving-view}

\code{copy_timestep} 把窗口中的一行复制成一个独立的
\code{MovingGreenTimeStep}。视图上的操作有两种：
\begin{itemize}
    \item 读取：\code{mblock(deltj)} 返回该行第 $\deltj$ 个块；
    \item 写回：把视图的结果拷回窗口的逻辑行。
\end{itemize}
写回时 \code{generation} 再次加一，任何在写回前取得的引用都应当
被视为过期。这一机制把“就地修改会破坏别名关系”的问题显式化：
视图是独立对象，写回是唯一合法的同步点。

\section{代数运算}
\label{sec:moving-algebra}

窗口对象与视图都支持逐元素线性组合：
\begin{align}
    \code{incr_step} &: X \leftarrow X + \alpha Y,\\
    \code{smul_step} &: X \leftarrow \alpha X,
\end{align}
其中 $\alpha$ 是复标量。它们只作用于视图的两个分量（推迟、小于），
不触碰 Matsubara 与左混合。窗口对象上的对应操作在内部对每一行
调用视图版本。

\section{乘作用}
\label{sec:moving-multiply}

时间相关的单粒子函数 $\bm{H}(t)$ 或 $\bm{f}(t)$ 以
\code{MovingTimeFunction} 表示。库提供两种乘作用：
\begin{align}
    \code{lmul_step} &: X(t,t') \leftarrow f(t)\,X(t,t'),\\
    \code{rmul_step} &: X(t,t') \leftarrow X(t,t')\,f(t'),
\end{align}
当 \code{adj} 为真时使用 $f^{\dagger}$。两者的实现都按行调用
块乘法，时间参数由相对指标决定：左乘取 $f$ 的槽 $\delti$，
右乘取槽 $\delti+\deltj$。

\section{时间外推}
\label{sec:extrapolation}

时间推进到一个新的领头步时，窗口需要一行“未来”的初值作为迭代
起点。库用最近 $k+1$ 行的外推给出：
\begin{equation}
    G(t_0-\delti h,\ \cdot) \approx
    \sum^{k}_{l=0} p_l\,
    G(t_0-(l+1)h,\ \cdot) .
    \label{eq:extrapolation}
\end{equation}
权重 $p_l$ 由 $k$ 阶外推多项式确定。实现要求
\begin{equation}
    t_c \ge 2k+1 ,
\end{equation}
因为外推需要 $k+1$ 行历史，而这些行本身在窗口末端要用到
$t_c-k$ 个相对偏移。靠近窗口末端时，索引要平移
\begin{equation}
    j_l =
    \begin{cases}
        j, & j < t_c-k,\\
        j-l-1, & j\ge t_c-k,
    \end{cases}
\end{equation}
以免同一行被重复使用。\code{extrapolate_timestep_moving}
接受阶数 $0..5$；阶数为 $0$ 时直接复制第 1 行。

\section{距离与检查}
\label{sec:moving-checks}

移动窗的距离有两种：
\begin{itemize}
    \item \code{distance_moving}：两个时间步视图之间的逐块
          Frobenius 距离之和；
    \item \code{distance_moving_at}：两个窗口在同一个绝对时间步上
          的对应行之间的距离，先按各自的 \code{tstp} 换算相对指标。
\end{itemize}
求解器入口还会检查窗口的兼容性：窗口长度、维数、统计符号与
领头时间步必须一致。几何检查由
\code{check_moving_solver} 完成，$\bm{H}$ 的检查由
\code{check_hamiltonian_moving} 完成。

\section{截断卷积}
\label{sec:moving-convolution}

移动窗上的卷积只访问窗口内的块。以推迟分量为例，
\begin{equation}
    C^{R}(n,m)
    =
    h\sum^{n}_{j=\max(m,\,n-t_c)}
      \Omm_{n-m,\,j-m}\,
      A^{R}_{n,j}\,f_j\,B^{R}_{j,m},
\end{equation}
求和下限被窗口截断：$j<n-t_c$ 的项被丢弃。内核对每个固定行
$n$ 只遍历窗口内的 $j$，因此单步代价为 $O(t_c)$，与总时间步数
无关。

\section{截断 Dyson}
\label{sec:moving-dyson}

\code{dyson_timestep_moving} 在窗口内求解推迟与小于分量。推迟
分量按列推进：每列 $m$ 的未知块只依赖窗口内的历史，因此系统
规模为 $O(t_c)\times O(t_c)$。小于分量的源项由窗口内的
$[\Sig^{R}\Gless+\Sig^{<}\Gadv]$ 累加得到，再逐列求解。

与完整解相比，截断解丢弃了两类项：超出窗口的卷积贡献，以及
窗口边界处的混合初始关联。第\ref{sec:moving-consistency}~节
给出它们何时可忽略。

\section{截断 VIE2}
\label{sec:moving-vie2}

\code{vie2_timestep_moving} 的结构与 Dyson 版本相同：推迟分量用
$k\times k$ 的启动系统与逐列推进，小于分量用 $(k+1)\times(k+1)$
的启动系统与逐列推进。核 $F$ 与源 $Q$ 都按移动窗存储。该接口
支持各阶求解，但为了保持逐位可复现，内部不做并行。

\section{截断气泡图}
\label{sec:moving-bubble}

气泡图也有移动窗版本：\code{bubble1_moving}、
\code{bubble2_moving} 及其厄米变体。它们的轨道指标约定与完整
版本相同（0-based），但只写入窗口内的行。厄米变体利用
\begin{equation}
    C^{R}(t,t') = -\big[C^{R}(t',t)\big]^{\dagger}
\end{equation}
从一半的数据重构另一半，减少约一半的乘法。

\section{截断观测量}
\label{sec:moving-observables}

截断能量积分以两个移动窗对象 $G$ 和 $D$ 为输入。记 $X_l$ 为
对象 $X$ 的相对索引 $(0,l)$ 块，则领头时间步的结果为
\begin{equation}
    E = h\,\mathrm{Im}\,\Tr\!\Big[\sum_{l=0}^{t_c}w_l
    \big((G_l^R)^\dagger D_l^<-(G_l^<)^\dagger D_l^R\big)\Big],
\end{equation}
这里 $h$ 是时间步长，$w_l$ 是 Gregory 权重。实现见 \code{kinetic_energy_trunc}。
密度矩阵和电流取自等时块，仍可能受到截断传播误差的影响。

\section{移动窗输入输出}
\label{sec:moving-io}

移动窗的文本格式是 \code{KB_MOVING_V1}，头部包含窗口长度、
初始化标志、领头时间步、统计符号与维数。数据按逻辑行顺序写出，
读取时重新建立环形映射。与完整轮廓格式一样，读取端逐行校验
列数、索引范围与数值有限性。

\section{一致性条件}
\label{sec:moving-consistency}

截断解与完整解在两种极限下一致：
\begin{enumerate}[label=(\arabic*)]
    \item 窗口覆盖全部历史，即 $t_0\le t_c$。此时没有项被丢弃，
          截断解与完整解逐位相同（除求和顺序可能不同）。
    \item 系统已经弛豫，$G(t,t')$ 对 $|t-t'|>t_c h$ 的依赖低于
          目标精度。此时丢弃的项在数值上可忽略。
\end{enumerate}
在两者之间，误差有两个来源：
\begin{itemize}
    \item \textbf{直接截断误差}：被丢弃的历史项本应贡献到当前步；
    \item \textbf{传播误差}：外推得到的初始行带有误差，它会在后续
          步中放大或衰减，取决于系统的稳定性。
\end{itemize}
第\ref{ch:validation}章给出与完整解对比的数值结果，包括不同窗口长度下的
误差随时间的演化。
